% !PRINTED 122
\weberpage{122}
%<<CONT
\[
\alpha_1 = \begin{pmatrix}1, & 0\\ 0, & 1\end{pmatrix}, \quad \alpha_2 = \begin{pmatrix}0, & 1\\ -1, & 0\end{pmatrix}, \quad \alpha_3 = \begin{pmatrix}1, & 1\\ -1, & 0\end{pmatrix},
\]
\[
\alpha_4 = \begin{pmatrix}1, & 1\\ 0, & 1\end{pmatrix}, \quad \alpha_5 = \begin{pmatrix}1, & 0\\ 1, & 1\end{pmatrix}, \quad \alpha_6 = \begin{pmatrix}0, & 1\\ -1, & 1\end{pmatrix},
\]
then $\alpha_1\mathfrak{A}, \alpha_2\mathfrak{A}, \alpha_3\mathfrak{A}, \alpha_4\mathfrak{A}, \alpha_5\mathfrak{A}, \alpha_6\mathfrak{A}$ are the cosets of $\mathfrak{A}$, and we have

\[
\mathfrak{L} = \alpha_1\mathfrak{A} + \alpha_2\mathfrak{A} + \alpha_3\mathfrak{A} + \alpha_4\mathfrak{A} + \alpha_5\mathfrak{A} + \alpha_6\mathfrak{A}
\]
and $\mathfrak{A}$ is a subgroup of $\mathfrak{L}$ of finite index 6:

\[
(\mathfrak{L}, \mathfrak{A}) = 6 \quad (\text{Bd.\ II}, \ \S\,2).
\]
The totality $\alpha_1\mathfrak{A} + \alpha_2\mathfrak{A}$ is likewise a group $\mathfrak{A}'$, and we have

\[
\mathfrak{L} = \alpha_1\mathfrak{A}' + \alpha_3\mathfrak{A}' + \alpha_4\mathfrak{A}'
\]
and $(\mathfrak{L}, \mathfrak{A}') = 3$.

\gesperrt{Just as all the linear transformations may be derived from the two fundamental transformations, so the transformations of the first class may be derived from repeated application of}

\[
\begin{pmatrix}1, & 0\\ 2, & 1\end{pmatrix}, \quad \begin{pmatrix}1, & 2\\ 0, & 1\end{pmatrix}, \quad \begin{pmatrix}-1, & 0\\ 0, & -1\end{pmatrix}
\]
\gesperrt{which can be proved by the method of \S\,30.}

\begin{center}
\textbf{\S\,37. \gesperrt{Representation of the $\sigma$-functions by $\vartheta$-functions.}}
\end{center}

In order to represent the function $\sigma(u)$ as a $\vartheta$-function, one can simply apply formula (6) of \S\,35 to the function

\[
t(u) = \vartheta_{11}\left(\frac{u}{\omega_1}, \frac{\omega_2}{\omega_1}\right)
\]
and thus set

\weq{1}{\sigma(u, \omega_1, \omega_2) = \omega_1\, e^{-\dfrac{u^2}{6\omega_1^2}\,\dfrac{\vartheta_{11}'''}{\vartheta_{11}'}}\; \frac{\vartheta_{11}\left(\dfrac{u}{\omega_1}, \dfrac{\omega_2}{\omega_1}\right)}{\vartheta_{11}'}}
But in consequence of the differential equation \S\,20, (4), we have:

\[
\vartheta_{11}''' = 4\pi i\, \frac{d\vartheta_{11}'}{d\omega}
\]
and hence, by the definition of the function $\eta(\omega)$ in \S\,34, (2):

\weq{2}{\frac{\vartheta_{11}'''}{\vartheta_{11}'} = 4\pi i\, \frac{d\log \vartheta_{11}'}{d\omega} = 12\pi i\, \frac{d\log \eta(\omega)}{d\omega}.}
